\chapter{Exercices supplémentaires : reconstruire les démonstrations}
Les exercices 1 à 20 du cours initial sont conservés. Les exercices suivants mettent l'accent sur les \emph{fondements} et la capacité à \emph{démontrer sans réciter}. Leur correction figure après les vingt premiers corrigés.

\exercice{21}{Fondements --- circonférence et $\pi$}
Expliquer pourquoi le quotient $C/D$ ne dépend pas du cercle. Calculer exactement la circonférence et l'aire d'un cercle de rayon $4\,\mathrm{cm}$. Préciser pourquoi $2\pi$ ne peut pas désigner la circonférence de tous les cercles.

\exercice{22}{Fondements --- homothétie}
Un cercle de rayon $2$ a une circonférence $C_1$ et une aire $A_1$. On double son rayon. Que deviennent la circonférence, le diamètre et l'aire ? Justifier les trois relations sans calculer $\pi$.

\exercice{23}{Fondements --- encadrement de $\pi$}
Pour le cercle unité, prouver que l'hexagone inscrit a un périmètre de $6$ et l'hexagone circonscrit de $4\sqrt3$. Déduire un encadrement exact de $\pi$.

\exercice{24}{Fondements --- radians}
Justifier l'égalité $180^\circ=\pi\,\mathrm{rad}$ en utilisant la définition du radian. Convertir $240^\circ$ en radians et calculer la longueur et l'aire d'un secteur de rayon $3\,\mathrm{cm}$ et d'angle $240^\circ$.

\exercice{25}{Seconde --- identité et domaine}
Démontrer $\sin^2x+\cos^2x=1$ à partir du cercle unité pour tout réel $x$. Si $\cos x=-3/5$ et $x$ appartient au deuxième quadrant, déterminer $\sin x$ et $\tan x$.

\exercice{26}{Seconde --- valeurs remarquables}
Sans utiliser de tableau, construire un triangle rectangle isocèle puis en déduire $\sin(\pi/4)$ et $\cos(\pi/4)$. Expliquer pourquoi $\cos(3\pi/4)=-\sqrt2/2$ et $\sin(3\pi/4)=\sqrt2/2$.

\exercice{27}{Première --- symétries}
Démontrer géométriquement $\sin(\pi-x)=\sin x$ et $\cos(\pi-x)=-\cos x$. En déduire exactement $\sin(5\pi/6)$ et $\cos(5\pi/6)$.

\exercice{28}{Première --- addition}
À partir des formules d'addition, démontrer $\sin(2x)=2\sin x\cos x$, puis $\cos(2x)=1-2\sin^2x$. En déduire $\sin^2x=(1-\cos2x)/2$.

\exercice{29}{Première --- tangente}
Démontrer $1+\tan^2x=1/\cos^2x$, en indiquant son domaine. Déduire $\tan(2x)=2\tan x/(1-\tan^2x)$ en précisant les conditions de validité.

\exercice{30}{Première --- équation par factorisation}
Résoudre dans $[0,2\pi[$ l'équation $\sin(3x)=\sin x$. Indiquer la factorisation et vérifier le nombre de solutions trouvées.

\exercice{31}{Première --- inéquation}
Résoudre dans $[0,2\pi]$ l'inéquation $\sin x\ge-\sqrt2/2$. Faire apparaître les angles frontières sur le cercle.

\exercice{32}{Première --- amplitude}
Démontrer que $2\sin x+2\cos x=2\sqrt2\sin(x+\pi/4)$. En déduire exactement les bornes des valeurs et tous les points de maximum dans $\mathbb R$.

\exercice{33}{Première --- triangles quelconques}
Dans un triangle, $b=3$, $c=4$ et $A=60^\circ$. Établir la loi des cosinus en choisissant des coordonnées convenables ; calculer ensuite $a$ et l'aire.

\exercice{34}{Première --- produit scalaire}
Soient $\vec u=(3,4)$ et $\vec v=(4,-3)$. Calculer leur produit scalaire, leurs normes et l'angle qu'ils forment. Expliquer le lien entre produit scalaire nul et perpendicularité.

\exercice{35}{Terminale --- limite}
À partir de l'inégalité géométrique $\sin h\le h\le\tan h$ pour $0<h<\pi/2$, démontrer $\lim_{h\to0}\sin h/h=1$. Expliquer pourquoi ce résultat nécessite de mesurer $h$ en radians.

\exercice{36}{Terminale --- dérivée et primitive}
Démontrer la formule $\bigl(\sin(3x)\bigr)'=3\cos(3x)$. Calculer une primitive de $\sin x\cos x$ de deux façons différentes et expliquer pourquoi elles peuvent différer d'une constante.

\exercice{37}{Terminale --- fonctions réciproques}
Calculer $\arcsin(\sin(7\pi/6))$, $\arccos(\cos(7\pi/6))$ et $\arctan(\tan(7\pi/6))$. Indiquer les intervalles de valeurs employés.

\exercice{38}{Terminale --- vibrations}
Une grandeur est décrite par $y(t)=4\cos(5\pi t-\pi/2)+1$, avec $t$ en secondes. Déterminer amplitude, période, fréquence, axe médian et valeurs extrêmes. Vérifier que $z(t)=y(t)-1$ satisfait $z''+25\pi^2z=0$.

\exercice{39}{Approfondissement --- formule de Moivre}
Démontrer $\cos(3x)=4\cos^3x-3\cos x$ en développant $(\cos x+i\sin x)^3$ et en identifiant les parties réelles.

\exercice{40}{Approfondissement --- corde et arc}
Deux points d'un cercle de rayon $R$ sont séparés par un angle au centre $\theta\in[0,\pi]$. Démontrer que la longueur $d$ de leur corde vaut $2R\sin(\theta/2)$ ; comparer $d$ à la longueur d'arc $R\theta$ et discuter l'égalité.
